\subsection{Generalization across molecular sizes}\label{sec:wide}
We evaluate the frozen source models on 64 further compositions, with 16 in each
of four size ranges from 17 to 64 atoms. All are absent from both processed OMol25
corpora and from the earlier generation panels. The models receive no additional
training. Metadata ranks select each composition after the same reference-validity
checks, with 32 draws per composition and model continuation.

\begin{table}[t]
\centering
\caption{Raw chemical-graph validity on the 64-composition panel. Each size range
contains 1,024 attempts per method across two continuations. Values are percentages;
the final column reports percentage-point differences.}
\begin{tabular}{lrrr}
\toprule
Atoms & Gaussian & Harmonic & Difference\\
\midrule
17--28 & 49.12 & 55.08 & +5.96\\
29--40 & 24.02 & 38.28 & +14.26\\
41--52 & 16.50 & 21.29 & +4.79\\
53--64 & 9.77 & 12.40 & +2.64\\
\midrule
All 64 compositions &24.85&31.76&+6.91\\
\bottomrule
\end{tabular}\label{tab:wide}
\end{table}

The harmonic source improves the pooled rate in every size range. Overall it
adds 6.91 percentage points, with a size-stratified composition-bootstrap 95\%
interval of [4.74,9.03]. The two continuations improve by 11.18 and 2.64 points,
respectively. Absolute validity decreases for larger molecules, while the source
advantage remains positive in the aggregate. These results extend the source
comparison beyond the original 17--28-atom confirmation range.
